\section{Inverse-power Riesz interactions}\label{sec:riesz}

The construction is not specific to the Coulomb exponent or to three
spatial dimensions. For an integer \(d\ge2\) and a real number \(s>0\),
define on \((\R^d)^3\)
\begin{equation}\label{eq:riesz-cost}
 c_s(x_1,x_2,x_3)=\sum_{1\le i<j\le3}|x_i-x_j|^{-s},
\end{equation}
with \(c_s=+\infty\) if two coordinates coincide. For a Borel probability
measure \(\mu\) on \(\R^d\), let \(\Pi(\mu,\mu,\mu)\) now denote the
couplings on \((\R^d)^3\), and write
\[
 E_s(\mu)=\inf_{\pi\in\Pi(\mu,\mu,\mu)}\int c_s\dd\pi.
\]

\begin{theorem}\label{thm:riesz}
For every integer \(d\ge2\) and every real \(s>0\), there is a
nonnegative density \(\rho\) on \(\R^d\) such that
\[
 \int_{\R^d}\rho(x)\dd x=1,\qquad
 \rho\in C_c^\infty(\R^d),\qquad
 \sqrt{\rho}\in C_c^\infty(\R^d),
\]
with the following properties. For \(\mu=\rho\dd x\), the value
\(E_s(\mu)\) is finite and attained by a coupling, but every pair of
Borel maps \(T_2,T_3:\R^d\to\R^d\) with
\((T_2)_\#\mu=(T_3)_\#\mu=\mu\) satisfies
\[
 \int_{\R^d}c_s\bigl(x,T_2(x),T_3(x)\bigr)\dd\mu(x)>E_s(\mu).
\]
The left-hand side may be infinite.
Nevertheless,
\[
 \inf_{\substack{T_2,T_3:\R^d\to\R^d\ \mathrm{Borel}\\
                 (T_2)_\#\mu=(T_3)_\#\mu=\mu}}
 \int c_s(x,T_2(x),T_3(x))\dd\mu(x)=E_s(\mu),
\]
and pairs with finite cost approach this infimum.
\end{theorem}

The density and all local construction constants may depend on
\((d,s)\). The physical Coulomb discussion in
Section~\ref{sec:introduction} refers to the specialization
\(d=3,s=1\); the interaction in \eqref{eq:riesz-cost} is the stated
inverse power, not a logarithmic kernel.

\begin{proof}
Fix \(d\ge2\) and \(s>0\), and put \(a_1=e_1\) and \(a_2=e_2\)
in \(\R^d\). The measure argument already applies in this dimension.
We verify the three properties of the cost that it needs: local
equality families given by diffeomorphisms, strict supporting
inequalities on all other component products, and a uniform
cost-oscillation bound on separated products.

\paragraph{Local branches.}
For \(h_s(b)=|b|^{-s}\), \(b\ne0\), the center jets are
\[
 c_s(0,a_k,-a_k)=2+2^{-s},\qquad
 \nabla_xc_s(0,a_k,-a_k)=0,\qquad
 \nabla_zc_s(0,a_k,-a_k)=s(1+2^{-s-1})a_k.
\]
As before, we keep the positive outer point \(y\) free and minimize
in the other two coordinates. The potentials in Lemma~\ref{lem:local}
are replaced by
\[
 u_0(x)=2+2^{-s}-\frac K2|x|^2,\qquad
 v_k(z)=s(1+2^{-s-1})a_k\cdot(z+a_k)-\frac K2|z+a_k|^2.
\]
The central potential is common to the two branches, and
\(F_k=c_s-u_0-v_k\) has vanishing \((x,z)\)-gradient at each center.
Direct differentiation gives
\[
 D^2h_s(b)=s|b|^{-s-2}
 \left((s+2)\frac{bb^{\mathsf T}}{|b|^2}-I_d\right).
\]
Its radial eigenvalue is \(s(s+1)|b|^{-s-2}\), and its tangential
eigenvalues are \(-s|b|^{-s-2}\), so it is invertible.

At the \(k\)-th center, define \(M_k\) and \(N_k\) as in
Section~\ref{sec:local}, with \(c_s\) in place of \(c\).
The two \(d\times d\) blocks of \(N_k\) are
\(-D^2h_s(-a_k)\) and \(-D^2h_s(-2a_k)\), hence are invertible,
while \(M_k\) is a symmetric \(2d\times2d\) matrix. For sufficiently
large \(K\), both \(M_k+KI_{2d}\) are positive definite and both
\(d\times d\) blocks of \((M_k+KI_{2d})^{-1}N_k\) are invertible, because
\[
 K(M_k+KI_{2d})^{-1}N_k
 =(I_{2d}+K^{-1}M_k)^{-1}N_k\longrightarrow N_k.
\]
One \(K\) works for the two branches. The implicit and inverse function
arguments of Lemma~\ref{lem:local} now give separate local
diffeomorphisms \(X_k,Z_k\) on \(d\)-dimensional neighborhoods.
The cost is smooth near the collision-free center triples; after
shrinking the product balls, the \((x,z)\)-Hessian of \(F_k\) remains positive
definite. Strict convexity on the product of the \(x\)- and \(z\)-balls
therefore gives the same local supporting inequality and exact equality
set, with \(w_k(a_k)=0\).

\paragraph{The supporting inequality and the marginal.}
The five center values of the potential are now
\(2+2^{-s},0,0,0,0\). A distinct-center triple containing zero but
not an opposite pair has cost
\[
 2+2^{-s/2}>2+2^{-s}.
\]
A distinct-center triple without zero has positive cost and potential
sum zero. These finitely many strict gaps persist on small products
by continuity. If two points lie in one ball of radius at most \(r\),
their pair cost is at least \((2r)^{-s}\). This exceeds \(3M\) for
sufficiently small \(r\), where \(M\) bounds the absolute values of
the preliminary potentials. Positivity of the other pair terms and
permutation symmetry of \(c_s\) therefore give the full supporting
inequality and exact contact set described in
Proposition~\ref{prop:certificate}, now in \(\R^d\).
The proof of Lemma~\ref{lem:common-ball} then also applies in \(\R^d\).

On that common ball \(B\), choose a nonnegative
\(g\in C_c^\infty(B)\) with \(\int_B g^2\dd x=1\), and form \(\nu\) and
\(\mu\) by \eqref{eq:mu} using the four resulting diffeomorphisms.
The change-of-variables construction in Section~\ref{sec:proof},
now with \(d\)-dimensional Jacobians, gives the same smooth compactly
supported density and square root: each image function is supported
compactly inside its open component, and the five supports are
disjoint. The same random choice of a branch and permutation gives
an optimal coupling \(\pi_0\) with the three marginals equal to
\(\mu\). The bounded supporting potential proves its cost is finite
and forces every optimal coupling onto the contact set. Applying
Lemma~\ref{lem:splitting} to the four branch maps then excludes
every marginal-preserving Borel pair, exactly as in
Section~\ref{sec:proof}.

\paragraph{Approximation by preserving maps.}
The five compact component sets have a positive mutual separation
\(\eta\), and \(\pi_0\) uses only products of three distinct components.
Repeat the cell partition and redistribution from the proof of
Proposition~\ref{prop:equal-infima} in \(\R^d\), using
Lemma~\ref{lem:borel-transport}. The first-coordinate marginals used
to split the cells are absolutely continuous by Fubini in every such
dimension. The redistribution gives a graph coupling
\(\pi_\delta\) of a preserving Borel pair with the same mass as
\(\pi_0\) on each product of cells of diameter at most \(\delta\).
Every active product still has all pair distances at least \(\eta\).
For two triples \(x,y\) in one such product, the scalar mean value
theorem gives
\[
 \left||x_\ell-x_r|^{-s}-|y_\ell-y_r|^{-s}\right|
 \le 2s\delta\,\eta^{-s-1},\qquad
 |c_s(x)-c_s(y)|\le6s\delta\,\eta^{-s-1}.
\]
Both costs are bounded by \(3\eta^{-s}\) on the active products.
Summing their oscillations against the common product masses yields
\[
 \left|\int c_s\dd\pi_\delta-\int c_s\dd\pi_0\right|
 \le6s\delta\,\eta^{-s-1}\longrightarrow0.
\]
The reverse inequality for the infima holds because every graph plan
is a coupling. This proves the final assertion. In particular, no
restriction relating \(s\) to \(d\) is needed: the optimizer and
these approximating plans stay away from all collisions.
\end{proof}
